\chapter{C ABI de \filepath{libgraphmatch}}
\label{sec:cabi}

\section{Headers}

\begin{itemize}
    \item \filepath{libgraphmatch/include/graphmatch/types.h}:
    tipos opacos (\code{gm_graph_t}, \code{gm_matcher_t},
    \code{gm_results_t}) y enumeraciones.
    \item \filepath{libgraphmatch/include/graphmatch/graphmatch.h}:
    funciones exportadas.
\end{itemize}

La librería se distribuye como objeto dinámico
(\filepath{libgraphmatch.so} en Linux,
\filepath{graphmatch.dll} en Windows) o estático
(\filepath{libgraphmatch.a}).

\section{Funciones de versión e introspección}

\begin{itemize}
    \item \code{uint32_t gm_abi_version(void)} — versión entera
    monotónica. Cliente valida contra su \code{expected_version}
    al \code{dlopen}; mismatch $\Rightarrow$ abort.
    \item \code{gm_simd_isa_t gm_simd_isa(void)} — ISA activa
    (\code{GM_ISA_SCALAR}, \code{GM_ISA_AVX2}, \code{GM_ISA_AVX512}).
\end{itemize}

\section{Construcción de grafo}

\begin{itemize}
    \item \code{gm_graph_t* gm_graph_new(void)} — crea grafo vacío.
    \item \code{void gm_graph_add_node(gm_graph_t* g, uint32_t id,
    uint8_t entity_type)} — agrega nodo.
    \item \code{void gm_graph_add_edge(gm_graph_t* g, uint32_t src,
    uint32_t dst, uint8_t relation_type, int64_t timestamp_ms)} —
    agrega arista.
    \item \code{uint32_t gm_graph_intern(gm_graph_t* g,
    const char* symbol)} — interna un símbolo, devuelve id.
    \item \code{const char* gm_graph_resolve(const gm_graph_t* g,
    uint32_t id)} — inverso del intern.
    \item \code{void gm_graph_finalize(gm_graph_t* g)} — finaliza
    la construcción (ordena adjacency lists, precomputa
    estadísticas). \textbf{Precondición}: sin nodos o aristas
    duplicados tras este punto.
    \item \code{uint64_t gm_graph_node_count(const gm_graph_t*)}
    — cuenta nodos.
    \item \code{uint64_t gm_graph_edge_count(const gm_graph_t*)}
    — cuenta aristas.
    \item \code{void gm_graph_free(gm_graph_t* g)} — libera.
\end{itemize}

\textbf{Ownership}: el caller posee el puntero y debe liberarlo
con \code{gm_graph_free}. La envoltura Rust
(\code{core/matcher-ffi::GmGraph}) hace esto automáticamente vía
RAII.

\section{Matcher}

\begin{itemize}
    \item \code{gm_matcher_t* gm_matcher_new(void)} — crea matcher.
    \item \code{void gm_matcher_set_start_type(gm_matcher_t* m,
    uint8_t entity_type)} — tipo del primer nodo de la hipótesis.
    \item \code{void gm_matcher_add_step(gm_matcher_t* m,
    uint8_t entity_type, uint8_t relation_type)} — siguiente paso.
    \item \code{void gm_matcher_set_time_window(gm_matcher_t* m,
    int64_t window_ms)} — ventana temporal entre pasos.
    \item \code{gm_results_t* gm_matcher_run(gm_matcher_t* m,
    const gm_graph_t* g)} — ejecuta. \textbf{Precondición}:
    \code{g} debe estar finalizado.
    \item \code{void gm_matcher_free(gm_matcher_t* m)} — libera.
\end{itemize}

\section{Resultados}

\begin{itemize}
    \item \code{uint32_t gm_results_count(const gm_results_t* r)}
    — número de matches.
    \item \code{float gm_results_get_score(const gm_results_t* r,
    uint32_t idx)} — score SIMD del match.
    \item \code{uint32_t* gm_results_get_path(const gm_results_t*
    r, uint32_t idx, uint32_t* out_len)} — path como arreglo de
    node IDs. El puntero pertenece a \code{r}; no liberar.
    \item \code{int64_t* gm_results_get_timestamps(
    const gm_results_t* r, uint32_t idx, uint32_t* out_len)} —
    timestamps de las aristas del match.
    \item \code{char* gm_results_to_json(const gm_results_t* r)} —
    serializa todos los matches como JSON. Ownership transferido
    al caller; liberar con \code{gm_free_string}.
    \item \code{void gm_results_free(gm_results_t* r)} — libera.
\end{itemize}

\section{Utilidades}

\begin{itemize}
    \item \code{void gm_free_string(char* s)} — libera strings
    emitidos por la librería (\code{gm_results_to_json}).
\end{itemize}

\section{Concurrencia}

\begin{itemize}
    \item \code{gm_graph_t} \textbf{después de \code{finalize}}: es
    read-only y seguro para acceso concurrente (múltiples matchers
    pueden correr sobre el mismo grafo).
    \item \code{gm_graph_t} \textbf{antes de \code{finalize}}: no
    thread-safe. Construir desde un solo hilo.
    \item \code{gm_matcher_t} y \code{gm_results_t}: propiedad
    exclusiva del hilo que los creó. Un match en curso no puede
    compartirse.
\end{itemize}

\section{Errores}

La ABI no usa códigos de error tradicionales. Contratos:

\begin{itemize}
    \item Funciones de construcción (\code{add_node}, \code{add_edge})
    son total: ignoran silenciosamente duplicados o IDs fuera de
    rango.
    \item \code{gm_matcher_run} devuelve \code{NULL} si el matcher
    está mal configurado (hipótesis vacía, grafo no finalizado).
    \item \code{gm_results_get_*} con \code{idx} fuera de rango
    devuelve punteros \code{NULL} y \code{*out_len = 0}.
\end{itemize}

El cliente Rust (\code{core/matcher-ffi}) traduce estos contratos
a \code{Result<T, MatcherError>} y al typestate de \code{GmGraph}.

\section{Protocolo de mismatch (ADR-002)}

Al cargar la librería, el cliente debe:

\begin{enumerate}
    \item Llamar \code{gm_abi_version()}.
    \item Comparar contra \code{EXPECTED_ABI_VERSION} compilada.
    \item Si no coinciden, abortar con mensaje claro (``shipped
    ABI V, binary expects V'').
\end{enumerate}

Esto previene que un binario se corra con una \filepath{.so}
incorrecta. No hay mecanismo para ``degradar funcionalidad'': la
ABI es o la esperada, o el binario no arranca.
